\documentclass[10pt,a4paper]{amsart}
\usepackage[margin=19mm]{geometry}
\usepackage{amsmath,amssymb,mathtools}
\usepackage[scr=rsfs,cal=euler]{mathalfa}
\usepackage{xfrac}
\usepackage[unicode,hidelinks,bookmarks=false]{hyperref}
\newcommand{\ie}{i.e.\ }
\newcommand{\cf}{cf.\ }
\newcommand{\resp}{respectively\ }
\newcommand{\Subsection}{\subsection}
\providecommand{\nobreakdash}{\nobreak\hbox{-}}
\setlength{\emergencystretch}{2em}
\multlinegap0pt
\usepackage{amssymb}
\usepackage[shortlabels]{enumitem}
\newtheorem{theorem}{Theorem}
\newcommand{\F}{\mathbb{F}_2}
\newcommand{\Shat}{\operatorname{Shat}}
\newcommand{\abs}[1]{\left|#1\right|}
\newcommand{\set}[1]{\left\{#1\right\}}
\title{Cofilling Shattering: A Syndrome-Support Hierarchy for Check Erasures}
\author{Joshua Steier}
\date{Source: arXiv:2607.17028v1 (CC BY 4.0)}
\begin{document}
\maketitle

\noindent\emph{Source and licence.} Cofilling Shattering: A Syndrome-Support Hierarchy for Check Erasures. Version arXiv:2607.17028v1 (CC BY 4.0), distributed by arXiv under the Creative Commons Attribution 4.0 International licence, http://creativecommons.org/licenses/by/4.0/, as recorded in the arXiv licence metadata for this exact version and shown on the abstract page https://arxiv.org/abs/2607.17028v1. Full licence notice: \texttt{licenses/arxiv-2607.17028-CC-BY-4.0.txt}.\par\smallskip

\noindent\textit{Editorial scope.} Counted unit: the article's theorem thm:check-basis-separation (source0.tex lines 1645--1671). The article's complete original proof of it is delivered with it (source0.tex lines 1673--1695, one \texttt{proof} environment of 23 lines). No mathematics is written, repaired or re-derived: the statement and the proof are the article's own lines, in the article's own notation, and no retained line is substituted or re-flowed.  No further source-local context is needed: every cross-reference of every retained line resolves inside the two delivered spans.  Nothing was omitted from any retained span, so the package carries no omission record.  No retained line contains a citation, so the wrapper carries no bibliography and no bibliography entry.  No retained line contains a figure environment, an image-inclusion command or drawing code, so no image is delivered and the package declares no asset dependency.  Editorial formatting only, in addition: an AMS \texttt{amsart} preamble that restates the article's own declarations and macros used by the retained lines, namely the environments \texttt{theorem} on separate counters, together with the standard packages those lines need.  The retained environments print in this document as Theorem 1 in order of appearance, whereas the article prints the same items with the numbering of its own full document.  The complete native source of the article as distributed in the arXiv e-print for this exact version is bundled unchanged as \texttt{sources/205525/source0.tex}.
\begin{theorem}[Pair-repetition code and extremal check-basis separation]
\label{thm:check-basis-separation}
Let
\[
 C_n=\set{(x,x):x\in\F^n}\subseteq\F^{2n},
 \qquad H_0=\begin{bmatrix}I_n&I_n\end{bmatrix}.
\]
Then $H_0$ is a full-rank parity-check matrix of the direct sum of $n$
binary repetition codes, and
\begin{equation}\label{eq:standard-pair-repetition}
 \Shat_{q,s}(H_0)=
 \begin{cases}
 \mathsf N_2(q,s),&\mathsf N_2(q,s)\le n,\\
 +\infty,&\mathsf N_2(q,s)>n.
 \end{cases}
\end{equation}
For every $q\ge1$ and $s\ge2$, put $n=\mathsf N_2(q,s)$.  There is
$T\in\operatorname{GL}(n,2)$ such that $H_1=TH_0$ is a row-equivalent
parity-check matrix of the same code and
\begin{equation}\label{eq:check-basis-gap}
 \Shat_{q,s}(H_1)=q,
 \qquad \Shat_{q,s}(H_0)=n.
\end{equation}
Thus the hierarchy is not invariant under change of check basis even when the
kernel code, domain and check lengths, rank, and image code are fixed.  For
$q\ge2$ the separation is genuinely collective.
\end{theorem}

\begin{proof}
For $y\in\F^n$, a preimage under $H_0$ is $(a,b)$ with $a+b=y$.
Coordinatewise $\abs a+\abs b\ge\abs y$, with equality at $(y,0)$, so
\begin{equation}\label{eq:H0-lambda}
 \lambda_{H_0}(y)=\abs y.
\end{equation}
A feasible $q$-space is therefore a binary code of distance at least $s$, and
its common support is its effective length.  This proves the lower bound in
\eqref{eq:standard-pair-repetition}; embedding a shortest
$[\mathsf N_2(q,s),q,d\ge s]$ code proves attainability.

Now set $n=\mathsf N_2(q,s)$ and choose a shortest binary $[n,q,d\ge s]$
code $V\le\F^n$.  Since $s\ge2$, $n>q$.  Choose
$T\in\operatorname{GL}(n,2)$ with
$T(V)=E_q:=\operatorname{span}\{e_1,\ldots,e_q\}$.  For $H_1=TH_0$,
\[
 \lambda_{H_1}(y)=\lambda_{H_0}(T^{-1}y)=\abs{T^{-1}y}.
\]
Hence every nonzero element of $E_q$ has localization at least $s$, so
$\Shat_{q,s}(H_1)\le q$.  A $q$-dimensional subspace needs at least $q$
support coordinates, proving equality.  Invertibility of $T$ gives
$\ker H_1=\ker H_0=C_n$.
\end{proof}

\end{document}
